\bookchapter{第六十五章：項目一——吾之記賬本（本地版）}

\chapterintro{迎汝入首個「正經之項」。此記賬本雖不連網，然五臟悉備：能記一筆之賬、能觀流水、能按月統計、能刪改。成之，則「數據 邏輯 界面」之全流程已遍歷矣。勿遽抄代碼——先思之審而後動手，此吾儕一貫之則也。}

\booksection{一、思其所需：先問五事}

\begin{center}
\begin{tabularx}{\linewidth}{|X|X|}
\hline
\textbf{要素} & \textbf{答案} \tabularnewline
\hline
Who（為誰而作） & 自用，家人朋友亦可 \tabularnewline
What（作何事） & 記收支、觀流水、按月彙總、按類統計 \tabularnewline
Why（何為） & 練手 + 明錢之所往 \tabularnewline
How（何以為之） & Swift + SwiftUI + SwiftData + MVVM，純本地而無後端 \tabularnewline
Constraint（約束） & 一週內成之；最低須支 iOS 17（SwiftData 之檻） \tabularnewline
\hline
\end{tabularx}
\end{center}

\begin{center}\small\color{gray}記賬本所需五要素\end{center}

繼而析為「步步可用」之小程：(1) 建 SwiftData 數據層（能貯一賬）；(2) 列表頁（能觀流水）；(3) 錄入彈窗（能增賬）；(4) 刪與編；(5) 月度彙總之卡。每成一程，App 皆「可用之物」也。

\booksection{二、數據層：使「賬」有其所}

先定一賬之形——字段足用即可，勿貪多：

\begin{lstlisting}[language=Swift]
import SwiftData

@Model
class Record {
    var type: String      // "income" 收入 | "expense" 支出
    var amount: Double    // 金額（元）
    var category: String  // 分類：餐飲、交通、薪俸……
    var note: String      // 備註（可選）
    var date: Date        // 記賬日期

    init(type: String, amount: Double, category: String, note: String = "", date: Date = Date()) {
        self.type = type
        self.amount = amount
        self.category = category
        self.note = note
        self.date = date
    }
}
\end{lstlisting}

\begin{tcbinfo}{小貼士：金額何以用 Double？}
嚴謹之金融 App 用 Decimal（精確十進制）以避浮點之差（0.1+0.2$\neq$0.3 之舊惑）。練手之項用 Double 足矣；他日若作真金融之品，宜易為 Decimal 或「以分為單位之 Int」——此無數先輩所踐之坑也。

\end{tcbinfo}

\booksection{三、MVVM：界面與邏輯「分家」}

徑使界面操數據庫乎？可行，然界面之碼將冗長可憎。MVVM 之構使各司其職：View（界面）惟司繪，ViewModel（邏輯之樞）司數據，Model（SwiftData）司存貯。

View（SwiftUI 界面）

ViewModel（@Observable 邏輯中樞）

Model（SwiftData 數據層）

\begin{center}\small\color{gray}圖 4：MVVM 三層之構——數據單向而流，各層互不越界\end{center}

\begin{lstlisting}[language=Swift]
import Observation

@Observable
class RecordViewModel {
    var records: [Record] = []
    private let context: ModelContext

    init(context: ModelContext) { self.context = context }

    func load() {
        // SwiftData 之 #Predicate 篩選 + 排序
        let descriptor = FetchDescriptor<Record>(
            sortBy: [SortDescriptor(\.date, order: .reverse)]
        )
        records = (try? context.fetch(descriptor)) ?? []
    }

    func add(type: String, amount: Double, category: String, note: String) {
        let r = Record(type: type, amount: amount, category: category, note: note)
        context.insert(r)
        try? context.save()
        load()                       // 重載，界面自新
    }

    func delete(_ record: Record) {
        context.delete(record)
        try? context.save()
        load()
    }
}
\end{lstlisting}

界面中 @State private var vm = RecordViewModel(...)，遂徑以 vm.records 渲列表——數據一變，界面自刷。

\booksection{四、界面：列表頁 + 錄入彈窗}

\begin{lstlisting}[language=Swift]
struct RecordListView: View {
    @State private var vm: RecordViewModel
    @State private var showAdd = false

    var body: some View {
        NavigationStack {
            List {
                Section("本月結餘") {          // 匯總卡片
                    SummaryRow(records: vm.records)
                }
                Section("流水") {
                    ForEach(vm.records) { r in
                        RecordRow(record: r)
                    }
                    .onDelete { idx in
                        idx.forEach { vm.delete(vm.records[$0]) }
                    }
                }
            }
            .navigationTitle("吾之記賬本")
            .toolbar {
                Button { showAdd = true } label: {
                    Image(systemName: "plus")
                }
            }
            .sheet(isPresented: $showAdd) {
                AddRecordView(vm: vm)          // 錄入彈窗
            }
        }
    }
}
\end{lstlisting}

錄入彈窗之要：金額輸入框之鍵盤用 .keyboardType(.decimalPad)（數字鍵盤），收入/支出用分段之控件（Picker），分類以菜單擇之，點「保存」則調 vm.add(...) 並 dismiss() 以闔窗。

\begin{tcbwarn}{按：輸入之驗勿懈}
金額輸入框必驗之：空值、負數、“abc” 之類非法輸入皆當攔（以 Double(text) 審之，不遂則返 nil）。App 之崩，其常因之一，即「視用戶之輸為馴良」。謹記：凡用戶所輸，皆當如敵而察之。

\end{tcbwarn}

\booksection{五、編、刪與彙總}

\begin{itemize}
  \item 刪：List 之 .onDelete 左滑而刪 + 確認彈窗（誤刪者，記賬類 App 用戶之所深憾也）；
  \item 編：點行入編輯彈窗，預填舊值（易一字較重輸為省）；
  \item 彙總：以 filter + reduce 算「本月收入 / 支出 / 結餘」，頂部之卡陳之：
\end{itemize}

\begin{lstlisting}[language=Swift]
func monthSummary(_ records: [Record]) -> (income: Double, expense: Double) {
    let cal = Calendar.current
    let monthRecords = records.filter {
        cal.isDate($0.date, equalTo: Date(), toGranularity: .month)
    }
    let income = monthRecords
        .filter { $0.type == "income" }
        .reduce(0) { $0 + $1.amount }
    let expense = monthRecords
        .filter { $0.type == "expense" }
        .reduce(0) { $0 + $1.amount }
    return (income, expense)
}
\end{lstlisting}

\booksection{六、驗收：三層之檢 + 二進階之挑戰}

項既成，按「功能 邊界 構架」三層驗之：能記否、能刪否、重啟 App 而數據存否（功能）；金額輸 0 或負數，崩否（邊界）；界面與數據有無相混、字段有無硬編（構架）。

基礎版驗過，則自加二道「進階之饌」：

\begin{enumerate}
  \item 分類統計：增「統計」標籤頁，以 Swift Charts 繪分類支出為餅圖/柱狀圖——蘋果官方圖表之庫，數行碼而圖成；
  \item iCloud 同步：SwiftData 開 iCloud 同步，惟須配數項設置——賬本於 iPhone/iPad 之間自同步，頓覺「高級」矣。
\end{enumerate}

\begin{tcbpractice}{來動手：本週成之}
按「數據層 ViewModel 列表頁 錄入彈窗 彙總卡片」之序而動手，日進一程。困則以報錯原文擲於搜索引擎或 AI（尚憶第七篇《AI編程實戰篇》乎？述問題時，明言「期望之行、實際之行、復現之步」）。成則以之記賬一週——自作之 App 自用之，成就之感將推汝前行。

\end{tcbpractice}

\begin{tcbsummary}{本章小結}
\begin{itemize}
  \item 所需五要素 析為「步步能用」之小程，而後動手
  \item 數據層：@Model 之模 + modelContainer 之器 + FetchDescriptor 之詢
  \item MVVM 分職：View 惟司繪，@Observable 之 ViewModel 司邏輯，SwiftData 司存貯
  \item 界面：List + Section + toolbar + sheet 錄入彈窗
  \item 三層驗收（功能/邊界/構架）+ 進階挑戰（Swift Charts、iCloud 同步）
\end{itemize}

\end{tcbsummary}
